\clearpage
\fancyhead[L]{\small CONNAISSANCE DES TEMS.}
\newcommand{\ObservedVenusFraction}{\hbox{\vrule width0pt height7.600000bp depth0pt\special{pdf:bcontent}\special{pdf:code q 0 g 0.205405 4.313514 m 0.000000 4.108108 l 0.000000 3.902703 l 0.205405 3.697297 l 0.821622 3.697297 l 1.027027 3.902703 l 1.232432 3.902703 l 1.437838 3.697297 l 1.437838 3.286486 l 1.232432 3.081081 l 1.027027 3.081081 l 0.821622 2.875676 l 0.821622 2.054054 l 1.232432 1.643243 l 1.232432 1.232432 l 1.027027 1.027027 l 1.027027 0.205405 l 1.232432 0.000000 l 1.437838 0.000000 l 1.643243 0.205405 l 1.848649 0.205405 l 2.054054 0.000000 l 2.464865 0.000000 l 2.670270 0.205405 l 2.875676 0.205405 l 3.697297 1.027027 l 3.697297 2.670270 l 3.081081 3.286486 l 3.081081 3.697297 l 3.286486 3.902703 l 3.697297 3.902703 l 3.902703 3.697297 l 4.724324 3.697297 l 4.929730 3.902703 l 4.929730 4.108108 l 4.724324 4.313514 l 3.286486 4.313514 l 3.081081 4.108108 l 1.027027 4.108108 l 0.821622 4.313514 l h 2.259459 7.600000 m 1.027027 6.367568 l 1.027027 6.162162 l 1.232432 5.956757 l 1.437838 5.956757 l 1.848649 5.545946 l 1.848649 5.135135 l 2.259459 4.724324 l 2.875676 4.724324 l 3.081081 4.929730 l 3.081081 5.135135 l 2.875676 5.340541 l 2.875676 7.394595 l 2.670270 7.600000 l h f* Q}\special{pdf:econtent}\kern4.929730bp}}

\hypertarget{native-586-lower}{}
\begin{center}\small Göttingische gelehrte Anzeigen. Issue 40. Pages 385..396. 11 March 1815.\\[3pt]\rule{.72\linewidth}{.3pt}\end{center}

\emph{Connaissance des tems, ou des mouvemens célestes, à l'usage des astronomes et des navigateurs pour
l'an} 1816; \emph{ publiée par le bureau des longitudes.} Paris. Published by the widow Courcier. November 1813. 360 pages in large octavo.

We pass over the astronomical calendar and the regular articles in silence and review only the \emph{supplements}, which had been almost entirely absent for several years and are again numerous and interesting in this volume. First comes a remarkable paper on comets by Count \textsc{Laplace}. Here this great geometer takes the side of the hypothesis that these wonderful celestial bodies arise from the condensation of the same matter that forms the nebulous stars and appears to be so generally distributed throughout space. On this view, comets would stand in the same merely accidental relation to our solar system as, in the opinion of some naturalists, \SourceLeaf{587}meteorites do to the Earth. If they come far enough within reach of our Sun for its attraction to predominate, they are compelled to describe elliptical or hyperbolic orbits around it. This readily explains the great variety, or rather the complete lack of any pattern, in the \emph{orientation} of the observed cometary orbits. One possible objection to this hypothesis seems to arise from the fact that among all the orbits calculated so far, already well over a hundred, we have yet to identify a single clearly hyperbolic one. For the kind of conic section depends on the comet's initial speed at the point where it enters the sphere in which the Sun's attraction predominates, and it seems just as capable of being a hyperbola as an ellipse. The aim of the paper is to prove by rigorous probability calculations that, among the orbits of all comets that come close enough for us to observe them, only an exceedingly small fraction are hyperbolas with an axis small enough for the observations to identify them as hyperbolas. For the ratio of the probability of such a hyperbola, whose semimajor axis does not exceed 100 radii of the Earth's orbit, to the probability of the remaining cases, namely a hyperbola with a larger axis, a parabola or an ellipse, \textsc{Laplace} finds the expression (I)
\[1 : \frac{\pi-2}{10}\BareRoot{\frac{r(r+200)}{2D}} - 1\]
where $\pi$ denotes half the circumference of a circle of radius 1, $D$ the maximum perihelion distance for comets that can still be observed, and $r$ the radius of the sphere in which the Sun's attraction predominates. Setting $r = 100000$, $D = 2$, assumptions that can reasonably be allowed, he finds that one could bet 56 to 1 that no recognizable hyperbola will occur among a hundred observed cometary orbits. This assumes that all perihelion distances from 0 to 2 are equally possible. If we assume that the probability falls more and more as perihelion distances increase, the above ratio becomes still larger. Under a very plausible hypothesis for the decreasing probability of increasing perihelion distances, \textsc{Laplace} calculates that one could bet 82 to one that no recognizable hyperbola will occur among a hundred observed orbits.

The importance of the subject and the name of the author justify dwelling a little longer on a paper of only a few pages. We meant the foregoing merely as a historical account of its contents and purpose. What we have yet to add is intended only for those of our mathematical readers who study the paper itself and have it before them. Much as we are inclined to give the hypothesis itself preference over all the others proposed for these enigmatic celestial bodies, especially \textsc{Lagrange}'s in the Connaissance des tems for 1814, we cannot help making two remarks about the mathematical part of the paper. First, it seems to us that \textsc{Laplace} is too hasty when he says on p. 216 `la plus petite valeur de $V$ est celle qui rend nulle la quantité renfermée sous le radical précédent.' The smallest value of $V$ is in fact 0, though the expression (II)
\[1{}^{\displaystyle -}\frac{\BareRoot{(1-\frac{D}{r})}}{rV}\cdot\BareRoot{(rrVV(1+\frac{D}{r})-2D)}\]
for the probability that the perihelion distance falls between 0 and $2D$ does not hold from $V=0$. It holds only from the value of $V$ that \textsc{Laplace} calls the smallest. For smaller values of $V$, that probability is not imaginary as expression (II) would make it. It is \SourceLeaf{588}$= 1$, that is, \emph{certainty}. In our view, \textsc{Laplace} was wrong to disregard these smaller values of $V$. We must therefore add $1\times\int dV$, taken from $V = 0$ to $V = \BareRoot{\dfrac{2D}{1+\dfrac{D}{r}}}$, that is, the quantity $\BareRoot{\dfrac{2D}{1+\dfrac{D}{r}}}$ itself, to the integral that \textsc{Laplace} expanded. If we otherwise kept \textsc{Laplace}'s expansions unchanged, this would give the following ratio in place of expression (I) above.
\[1 : \frac{\pi}{10}\BareRoot{\frac{r(r+200)}{2D}} - 1\]
Thus one could actually bet 157 to 1 in place of 56, which would be still more favorable to the hypothesis. But we have a second remark to make, which produces a very different result. On p.\,217 \textsc{Laplace} expands the integral expression (b) that he found on p.\,216 into a series and obtains (III)
\[\frac{(\pi-2)\BareRoot{2D}}{2r} - \frac{D}{ir\BareRoot{r}}\]
which, according to our first remark, would have to be replaced by
\[
\frac{\pi\BareRoot{2D}}{2r} - \frac{D}{ir\BareRoot{r}}
\]
These expressions therefore approach finite limits as $i$ grows to infinity. This cannot be, as one can readily see from the exact integral or even from the differential before integration. Indeed, on carrying the expansion further, one finds that (IV) must also be added.
\[+\,\frac{iDD}{2rr\BareRoot{r}}\]
By disregarding this part, \textsc{Laplace} was able to make his final result independent of $U$ by treating that quantity as infinite (an assumption \textsc{Laplace} in fact makes tacitly when he takes $i$ to be infinitely large on p.\,217). But suppose one still allowed the same assumption that every value of the initial speeds is equally possible without restriction --- and took term (IV) into account. The final result would then be entirely different. It would be infinitely unlikely for any ellipse, parabola, or hyperbola approaching a parabola to occur among a finite number of observed cometary orbits. In all probability they would all be hyperbolas indistinguishable from straight lines. But the assumption itself is clearly inadmissible, since it would leave every finite speed with infinitely little probability. One \emph{must} assume a finite bound for $U$. Since the final result then depends entirely on $U$ and we have no grounds for settling its value, the problem is in fact insoluble. Yet even though probability theory cannot provide a decisive proof \emph{for} the hypothesis, for the very reason that we do not know the bound $U$, it also decides absolutely nothing against the hypothesis. The calculation itself shows that as long as one does not assume \emph{very} large values for $U$, recognizably hyperbolic orbits must always remain very rare. From a rough estimate we find that one must assume at least $U = 1083$ (the unit is the mean speed of the Earth in its orbit) to be justified in expecting to find one orbit that is at least clearly hyperbolic among a hundred cometary orbits.

\SourceLeaf{589}We may be briefer in reviewing the other papers. In the next one the learned editor of the C. d. T. asks whether \textsc{Hipparchus} made observations in Alexandria and makes it seem likely that this greatest astronomer of antiquity either was never in Alexandria or at most stayed there briefly. --- By the same author, on the azimuth observations made by \textsc{Tycho} at Uranienburg (which are admittedly rather inaccurate but date only from the early part of his stay on the island of Hveen and give us no grounds for doubting that his meridian instruments were properly set up). --- Remarkable extracts from the manuscript of \textsc{Ptolemy}'s Optics held in the Paris library. --- An extract from \textsc{Don Joseph Rodriguez}'s paper on the English measurement of a degree, with some remarks. Mr. \textsc{Rodriguez} sought to show that the partial discrepancies found in the various measurements of degrees do not yet entitle us to regard an irregular shape of the Earth as established. Mr. \textsc{Delambre} rightly observes in reply that the contrary has not been proved either. But what this astronomer adds about \textsc{Mechain}'s observations in Barcellona and Montjouy seems to us in no way to undermine the conclusion from those observations, namely that the elevations of the pole determined with \textsc{Lenoir}'s circles leave uncertainties of several seconds. To assess properly these remarkable observations, which have given rise to much debate back and forth, one should really compare the results as follows. The meridian arc between Montjouy and Barcellona, which the geodetic measurements gave $=59''53$, is obtained from the astronomical observations
\begin{center}
\fontsize{9}{11}\selectfont\setlength{\tabcolsep}{3pt}\begin{tabular}{@{}l c r@{}}
from $\beta$ Geminorum & $.\;.\;.\;.\;.\;.\;.\;.\;.\;.\;.\;.\;.\;.\;.\;.$ & $57''20$\\
$\zeta$ Ursae Majoris upper culmination & $.\;.\;.\;.\;.\;.\;.\;.$ & $61.59$\\
$\zeta$ Ursae Majoris lower culmination & $.\;.\;.\;.\;.\;.\;.\;.$ & $64.66$\\
$\beta$ Ursae Minoris upper culmination & $.\;.\;.\;.\;.\;.\;.\;.$ & $62.75$\\
$\beta$ Ursae Minoris lower culmination & $.\;.\;.\;.\;.\;.\;.\;.$ & $63.39$\\
$\alpha$ Ursae Minoris upper culmination & $.\;.\;.\;.\;.\;.\;.\;.$ & $62.42$\\
$\alpha$ Ursae Minoris lower culmination & $.\;.\;.\;.\;.\;.\;.\;.$ & $63.12$\\
\end{tabular}
\end{center}
Thus, independently of all local attractions, we find differences of up to $7''46$ between the results from different stars. These differences can be attributed only to the observation and the instrument. --- Some further remarks on Mr. \textsc{Olbers}'s familiar parallax formulas, by Mr. \textsc{Delambre}, and by the same author a notice of his Abrégé de l'astronomie and of Halma's translation of the Almagest. --- Improved elements of Jupiter's motions by Mr. \textsc{Bouvard}, from which we may expect new tables soon. --- Various remarks on measuring heights with the barometer, by Mr. \textsc{Prony}. A new method described here for increasing the accuracy of the readings is worth adopting. In several series of observations cited here, the extreme results for the difference in height differ by only 2 to 3 metres on average. --- On the pendulum experiments conducted by Spanish navigators at many points on the Earth, by Mr. \textsc{Mathieu}. --- On Mr. \textsc{Burckhardt}'s new lunar tables. We find here what we had wished for in our review of these excellent tables, namely the complete set of individual results of the comparison with the observations. It seems remarkable to us that negative errors in latitude noticeably predominate. We take this to mean (though it is not explicitly stated) that on average the tables put the latitude $2''29$ farther south than the observations do. This suggests the need for a small correction to the elements used in the calculation. But it will be hard to decide how much of that may be due to the parallax constant, the \SourceLeaf{590}collimation error, the refraction, or the elevation of the pole. --- On the masses of the planets, by the same author. In preparation for constructing new tables for Mars, Venus and Mercury, which we now see we may expect from this excellent astronomer, he has reexamined the elements of the Earth's motion using nearly 4000 observations. We have here the principal results of that work. It is remarkable that this implies a reduction of about $\ObservedVenusFraction$ in the mass of Venus, whereas Mr. \textsc{von Lindenau} found from the motions of Mercury an increase of about the same size. --- Remarks on multiplying circles, by Mr. \textsc{Arago}. All astronomers know the objections that Mr. \textsc{von Zach} has raised against circles with a fixed spirit level on the axis. Mr. \textsc{Arago} seeks to dispel them in this paper. He says that theory does seem to support the possibility that the operations involved in the even-numbered observations might slightly alter the connection between the circle and the column in a way the level cannot show. But the instrument makers themselves must be consulted as arbiters, and they do not believe any appreciable displacement can occur provided the clamp is the right size and the alidade moves smoothly along the limb. (The concern is less about the clamp than about the adjusting screw, which after all is the sole connection between the circle and the column. It seems to us that \emph{this} doubt could be cleared up if, in some series of observations of one and the same star, the telescope were always brought back to the star through the zenith as usual for the even-numbered observations, while in other series it were brought back to the star through the nadir by a movement in the opposite direction. If moving the telescope really caused an appreciable displacement of the circle itself, one would have to expect the results of the first series to differ from those of the others consistently in the same direction.) In his view, the experiences Mr. \textsc{von Zach} cites against circles with a fixed level are inconclusive here, because they come from other observers, whereas the observations with circles fitted with a movable level, which agree with one another so admirably, were made by Mr. \textsc{von Zach} himself. Mr. \textsc{Arago} also defends the method of observation that Mr. \textsc{von Zach} disapproves of. In the even-numbered observations the column itself is left as it is, and one merely records each time the position of the bubble against the scale divisions and corrects the final result accordingly. Mr. \textsc{Arago} examined the value of the divisions of the level's scale on the three-foot \textsc{Reichenbach} circle, which the Paris Observatory received as a gift from Count \textsc{Laplace}, on two occasions at very different temperatures and obtained results that agreed almost exactly. Still, this leaves \emph{the} doubt whether one is justified. in treating the inside of the level tube as curved so uniformly throughout that one may assign exactly equal values to \emph{each individual} division of the scale. If one considers how extraordinarily small a departure from regular curvature is enough here to produce differences of several seconds, one will find some grounds for doubting that even the most skilful instrument maker can avoid such departures altogether. In any case, until one is completely certain about this, we would think it safer to adjust the column each time for the even-numbered observation so as to bring the air bubble back to exactly the position it occupied in the preceding odd-numbered observation.

Mr. \textsc{Arago} also turns to Mr. \textsc{von Zach}'s remarkable claim that, however well observations with repeating circles agree with one another, they all still leave an uncertainty of several seconds in the absolute accuracy of the results, since \emph{different} circles give different results. Since this famous observer himself has declared the source of these differences to be still a riddle, \SourceLeaf{591}another astronomer, who is unaware even of some of the circumstances, would be still less able to settle the matter. Mr. \textsc{Arago} does, however, give some similar experiences of his own. He believes he finds their source in a circumstance on which we have a few more things to say. From 422 observations of the Pole Star with a \textsc{Fortin} circle with a fixed spirit level, comprising four series that agreed admirably with one another, the latitude of Formentera came out as $38^\circ\ 39'\ 56''05$. The objective was then shifted slightly to obtain a clearer view. Seven further series, again in excellent agreement with one another, gave $38^\circ\ 39'\ 53''91$ from 570 observations. Similar differences later appeared with the same instrument in Paris, where the elevation of the pole varied between $48^\circ\ 50'\ 11''$ and $48^\circ\ 50'\ 15''$ depending on the setting of the objective, although any particular setting always gave results that agreed extremely closely. The result was always closely tied to the setting of the objective, though in different ways for different observers. Mr. \textsc{Arago} concludes that these anomalies arise solely from the way the center of the star is estimated. Telescopes always show the image of the star much too large, and Mr. \textsc{Arago} doubts whether the center of the apparent image is also the center of the true one. Each observer, he says, has an individual way of judging the center, and that would account for the phenomenon very naturally. We must confess that we cannot accept this explanation. Even if we were to grant that the two centers really differ as described, it is clear that if the telescope were turned halfway about its axis, the center previously above the other would now be below it. That is exactly what must happen with the method of observing with repeating circles. The telescope takes opposite positions in the even- and odd-numbered observations, so the pointing errors must cancel. Even if the telescope produced an image of quite irregular shape, the particular point at which one aimed would make no difference at all, provided one always chose one and the same point in odd- and even-numbered observations. If we may take the reported fact as fully established, one might almost be inclined to suspect that shifting the objective moved the crosshairs outside its focus, so that they had parallax. In that case it would of course matter how the observer was accustomed to holding his eye in front of the eyepiece. It is striking in any case that the telescope was designed so that the objective had to be moved to obtain a clear view. The paper also gives interesting results for the elevation of the pole obtained from zenith distances with the three-foot \textsc{Reichenbach} circle mentioned above. One might only have wished to see the elements of the calculation given, so that these observations could also help determine both the declination and the right ascension of the Pole Star. Perhaps, however, we shall receive these observations in more detail in the next volume of the Connoissance des tems, for which an illustration of this excellent instrument is also promised.

